\documentclass[final]{beamer}

\usepackage[size=custom,width=182.88,height=91.44,scale=1.4]{beamerposter} % 72x36 inches in cm
\usepackage{graphicx}
\usepackage{booktabs}
\usepackage{tikz}
\usepackage{xcolor}
\usepackage{url}
\usepackage{hyperref}
\usepackage{multicol}
\usepackage{enumitem}
\usepackage[T1]{fontenc}
\usepackage{lmodern}
\usepackage{tcolorbox}
\tcbuselibrary{skins}

% ── Colour palette ──────────────────────────────────────────────────────────
\definecolor{teal}{RGB}{32,94,106}        % deep teal — header & block titles
\definecolor{tealdark}{RGB}{20,64,74}     % darker teal for header bg
\definecolor{teallight}{RGB}{232,244,246} % very light teal — block bodies
\definecolor{gold}{RGB}{220,160,30}       % warm gold — accent squares & arrows
\definecolor{golddark}{RGB}{180,120,10}   % darker gold for contrast
\definecolor{slate}{RGB}{90,100,110}      % neutral slate for secondary text
\definecolor{softblue}{RGB}{100,145,175}  % soft blue — architecture boxes

% ── Beamer theme ────────────────────────────────────────────────────────────
\usetheme{default}
\setbeamercolor{block title}{fg=white,bg=teal}
\setbeamercolor{block body}{fg=black,bg=teallight}
\setbeamercolor{block title alerted}{fg=white,bg=softblue}
\setbeamercolor{block body alerted}{fg=black,bg=white}
\setbeamerfont{block title}{size=\large,series=\bfseries}
\setbeamertemplate{navigation symbols}{}

% ── Helper: numbered circle ──────────────────────────────────────────────
\newcommand{\step}[1]{%
  \tikz[baseline=(x.base)]{
    \node[circle,draw=teal,fill=teal,text=white,
          inner sep=4pt,font=\bfseries\large] (x) {#1};
  }%
}

\begin{document}
\begin{frame}[t]{}

%% ═══════════════════════════════════════════════════════════════════════════
%%  HEADER
%% ═══════════════════════════════════════════════════════════════════════════
\begin{tikzpicture}[remember picture,overlay]
  \fill[tealdark] (current page.north west) rectangle
               ([yshift=-9cm]current page.north east);
\end{tikzpicture}

\vspace{-1cm}
\begin{columns}[T]
  \begin{column}{0.78\linewidth}
    {\color{white}\fontsize{54}{60}\selectfont\bfseries
      \raisebox{0.15em}{\textcolor{gold}{$\blacksquare$}}~BookLens:
      Human-Centered Book Dataset Exploration}\\[0.5em]
    {\color{white}\Large Joshua Keith \quad University of Utah \quad
      \texttt{u1423033@utah.edu}}\\[0.6em]
    {\color{white}\LARGE
      \href{https://youtu.be/placeholder}{\underline{Demo Video}}
      \qquad\qquad
      \href{https://github.com/joshkeith1/BookLens}{\underline{github.com/joshkeith1/BookLens}}}
  \end{column}
  \begin{column}{0.18\linewidth}
    \raggedleft
    % Replace with actual UofU logo if available:
    % \includegraphics[width=6cm]{uu_logo.png}
    \color{white}\Large\bfseries University\\of Utah
  \end{column}
\end{columns}

\vspace{1.2cm}

%% ═══════════════════════════════════════════════════════════════════════════
%%  THREE COLUMNS
%% ═══════════════════════════════════════════════════════════════════════════
\begin{columns}[T,totalwidth=\linewidth]

%% ── LEFT COLUMN ─────────────────────────────────────────────────────────
\begin{column}{0.30\linewidth}

  \begin{block}{Book data is inaccessible to non-experts}
    Readers generate rich personal data — hundreds of books, ratings, genres,
    and formats — but have no practical way to explore it meaningfully.

    \vspace{0.5em}
    \begin{tcolorbox}[colback=white,colframe=gold,arc=3pt,boxrule=1.5pt]
      \textit{``What genres have I read most this decade? Do longer books get
      lower ratings? How has my reading pace changed year over year?''}

      \vspace{0.4em}
      These are natural questions. Answering them today requires SQL or data
      science skills that most readers simply do not have.
    \end{tcolorbox}
  \end{block}

  \vspace{0.6em}

  \begin{block}{Why is this hard to solve?}
    \begin{itemize}[leftmargin=1.2em,itemsep=0.4em]
      \item \textbf{Heterogeneous schemas} — column names and types vary
            widely across datasets (\texttt{user\_rating} vs.\
            \texttt{average\_rating}, \texttt{name} vs.\ \texttt{title})
      \item \textbf{Adaptive chart selection} — appropriate visualizations
            depend on which columns are present; a system must decide
            dynamically what to show
      \item \textbf{Narrative gap} — raw statistics are opaque without
            plain-language interpretation
      \item \textbf{Deployment friction} — tools requiring installation
            are a barrier before a user even begins
    \end{itemize}
  \end{block}

  \vspace{0.6em}

  \begin{block}{Applications \& Users}
    \begin{itemize}[leftmargin=1.2em,itemsep=0.4em]
      \item \textbf{Casual readers} tracking personal libraries (Goodreads
            exports, spreadsheets)
      \item \textbf{Students} exploring public book datasets for coursework
      \item \textbf{Educators} demonstrating data exploration concepts
            without coding prerequisites
      \item \textbf{Librarians} analyzing collection metadata and usage trends
    \end{itemize}
  \end{block}

  \vspace{0.6em}

  \begin{alertblock}{Impact}
    By removing every technical barrier between a user and their data,
    BookLens makes data-driven reading insights accessible to anyone — no SQL,
    no configuration, no installation required.
  \end{alertblock}

\end{column}

\hfill

%% ── MIDDLE COLUMN ───────────────────────────────────────────────────────
\begin{column}{0.36\linewidth}

  \begin{block}{BookLens enables zero-code data exploration}
    \centering
    % Replace with actual screenshot:
    \fbox{\parbox{0.92\linewidth}{
      \centering\vspace{3cm}
      \textit{[Insert annotated BookLens screenshot here.\\
      Recommended: sidebar open on left, Visualizations tab\\
      showing Genre Over Time + Reading Format Over Time.\\
      Add circled callouts \step{1}–\step{5} on the image.]}
      \vspace{3cm}
    }}
  \end{block}

  \vspace{0.6em}

  \begin{block}{BookLens supports human-in-the-loop exploration}
    \begin{itemize}[leftmargin=1.2em,itemsep=0.3em]
      \item Select from curated datasets or upload any personal CSV
      \item Apply sidebar filters interactively (genre, year, rating,
            keyword search)
      \item Switch datasets mid-session; filter panel adapts to the new schema
    \end{itemize}
  \end{block}

  \vspace{0.6em}

  \begin{block}{Demonstration Walkthrough}
    \vspace{0.3em}
    \step{1}~\textbf{Upload or select a dataset} from the sidebar —
    BookLens auto-detects columns and normalizes the schema with no
    user input required.

    \vspace{0.5em}
    \step{2}~\textbf{Overview \& Profiling} — review row counts,
    missing-value summaries, column types, and value distributions
    at a glance.

    \vspace{0.5em}
    \step{3}~\textbf{Adaptive Visualizations} — charts appear only for
    columns that exist: genre trends, reading pace, format over time,
    top authors, rating distributions, and more.

    \vspace{0.5em}
    \step{4}~\textbf{Cross-dataset comparison} — switch to Goodreads
    or Amazon bestsellers; the filter panel updates automatically
    to match the new schema.

    \vspace{0.5em}
    \step{5}~\textbf{AI Narrative Insights} — one click sends summary
    statistics to Claude (via OpenRouter); plain-language analysis
    streams back instantly in the browser.
  \end{block}

\end{column}

\hfill

%% ── RIGHT COLUMN ────────────────────────────────────────────────────────
\begin{column}{0.30\linewidth}

  \begin{block}{BookLens in a nutshell}

    \textbf{System Architecture}

    \vspace{0.5em}
    \newcommand{\archbox}[1]{%
      \begin{tcolorbox}[colback=softblue!10,colframe=softblue,
                        arc=4pt,boxrule=1.2pt,halign=center,
                        top=4pt,bottom=4pt]
        \bfseries #1
      \end{tcolorbox}%
    }
    \newcommand{\archarrow}{%
      \begin{center}
        {\color{gold}\Large$\downarrow$}
      \end{center}
      \vspace{-0.4em}
    }
    \archbox{CSV Input\\(upload or built-in)}
    \archarrow
    \archbox{Data Loader\\Schema Normalization}
    \archarrow
    \archbox{Filter Engine\\(sidebar widgets)}
    \archarrow
    \archbox{Analytics Engine\\(5 tab renderers)}
    \archarrow
    \archbox{AI Insights Layer\\Claude via OpenRouter}

    \vspace{0.8em}
    \textbf{Three Core Contributions}
    \vspace{0.3em}

    \begin{tcolorbox}[colback=softblue!10,colframe=softblue,
                      arc=3pt,boxrule=1pt,left=6pt]
      \textbf{\textcircled{1} Schema Normalization Pipeline}\\
      Maps heterogeneous CSV conventions to a unified schema
      automatically — no manual column mapping needed.
    \end{tcolorbox}
    \vspace{0.3em}
    \begin{tcolorbox}[colback=softblue!10,colframe=softblue,
                      arc=3pt,boxrule=1pt,left=6pt]
      \textbf{\textcircled{2} Adaptive Visualization Engine}\\
      Inspects available columns and renders only meaningful
      charts from a catalog of 10 chart types.
    \end{tcolorbox}
    \vspace{0.3em}
    \begin{tcolorbox}[colback=softblue!10,colframe=softblue,
                      arc=3pt,boxrule=1pt,left=6pt]
      \textbf{\textcircled{3} AI Narrative Layer}\\
      Computes summary statistics and streams plain-language
      insights via Claude Sonnet (OpenRouter API).
    \end{tcolorbox}
  \end{block}

  \vspace{0.6em}

  \begin{block}{Datasets \& Performance}
    \begin{tabular}{lr}
      \toprule
      \textbf{Dataset} & \textbf{Size} \\
      \midrule
      Personal Library  & 381 books \\
      Goodreads         & $\sim$11,000 books \\
      Amazon Top 50 (2009–2019) & 550 entries \\
      Book-Crossings    & 271,000 books \\
      \bottomrule
    \end{tabular}

    \vspace{0.5em}
    Load time for 271K rows: \textbf{$<$ 3 seconds}\\
    Filter operations: \textbf{instantaneous}\\
    AI insight accuracy: \textbf{validated across all 4 datasets}
  \end{block}

  \vspace{0.6em}

  \begin{block}{References}
    \small
    [1] Kandel et al. \textit{Wrangler}. CHI 2011.\\[0.2em]
    [2] Luo et al. \textit{DeepEye}. ICDE 2018.\\[0.2em]
    [3] Vartak et al. \textit{SeeDB}. VLDB 2015.\\[0.2em]
    [4] Wongsuphasawat et al. \textit{Voyager}. IEEE TVCG 2016.\\[0.2em]
    [5] Zhong et al. \textit{Seq2SQL}. arXiv 2017.
  \end{block}

\end{column}

\end{columns}

\end{frame}
\end{document}
